% Generated draft: requires mathematical and source review.
\documentclass[11pt]{article}
\usepackage[margin=25mm]{geometry}
\usepackage{fontspec}
\setmainfont{Latin Modern Roman}
\usepackage{amsmath,amssymb,mathtools,unicode-math}
\setmathfont{Latin Modern Math}
\usepackage{enumitem,needspace,xurl,hyperref}
\setlength{\parindent}{0pt}
\setlength{\parskip}{5pt}
\setlength{\emergencystretch}{3em}
\newcommand{\N}{\mathbb{N}}
\newcommand{\Z}{\mathbb{Z}}
\newcommand{\Q}{\mathbb{Q}}
\newcommand{\R}{\mathbb{R}}
\newcommand{\C}{\mathbb{C}}
\newcommand{\F}{\mathbb{F}}
\newcommand{\A}{\mathbb{A}}
\newcommand{\PP}{\mathbb{P}}
\newcommand{\E}{\mathbb{E}}
\newcommand{\eps}{\varepsilon}
\newcommand{\dd}{\,\mathrm{d}}
\begin{document}
\section*{UnsolvedMath ID: 30006623}
\textbf{Source problem code:} catalogue entry
\subsection{Definitions and mathematical statement}
A metric space $(X,d)$ is median if, for every $x,y,z\in X$, the intersection
\[
I(x,y)\cap I(y,z)\cap I(z,x),\qquad I(x,y)=\{w\in X:d(x,w)+d(w,y)=d(x,y)\},
\]
consists of exactly one point, denoted $m(x,y,z)$. A subset is convex if it contains $I(x,y)$ whenever it contains $x$ and $y$. A halfspace is a nonempty proper subset whose two complementary sides are convex. Two halfspaces $H,K$ are transverse if all four intersections of their complementary sides are nonempty. The rank of $X$ is the supremum of the sizes of finite families of pairwise transverse halfspaces; finite rank means that this supremum is finite.

A group $G$ is word-hyperbolic if it has a finite generating set for which its Cayley graph has $\delta$-thin geodesic triangles for some $\delta\geq 0$. Interpret the actions below as actions by isometries, subject to the source qualification below. An action is cocompact if $X=G K$ for some compact subset $K\subseteq X$. Two possible meanings of properness require distinction:
\begin{itemize}
\item compact-set properness: for every compact $K\subseteq X$, the set $\{g\in G:gK\cap K\ne\varnothing\}$ is finite;
\item metric properness: for every $x\in X$ and $R\geq 0$, the set $\{g\in G:d(x,gx)\leq R\}$ is finite.
\end{itemize}
The supplied record does not specify which meaning is intended, or assert that $X$ is a proper metric space. These conventions must therefore be checked before selecting a definitive formulation.

Use the usual geometric interpretation of \emph{cocompactly cubulable}: admitting a proper, cocompact action by cubical automorphisms on a CAT(0) cube complex. A group $G$ is residually finite if
\[
\forall g\in G\setminus\{1\}\ \exists\text{ a finite group }F\ \exists\phi:G\to F\text{ a homomorphism such that }\phi(g)\ne 1.
\]
Under the unrestricted-field interpretation, $G$ is linear if there exist a field $k$, an integer $n\geq 1$, and an injective homomorphism $G\hookrightarrow\operatorname{GL}_n(k)$. The record specifies neither a field nor a characteristic.

\textbf{Questions in the supplied curated formulation.} For every word-hyperbolic group $G$ admitting a proper, cocompact isometric action on a finite-rank median space $X$, must each of the following conclusions hold?
\begin{enumerate}
\item $G$ is cocompactly cubulable.
\item $G$ is residually finite.
\item $G$ is linear.
\end{enumerate}
Here properness remains subject to the explicit unresolved convention above. No completeness, geodesicity, local compactness, or prescribed rank bound for $X$ is supplied.

\textbf{Source qualifications.} The displayed questions follow the supplied statement and clean statement. The original excerpt contains only the cubulability and residual-finiteness questions, with their hypotheses missing; it does not contain the linearity question. The supplied literature assessment discusses residual finiteness for an action on a median space without explicitly retaining the finite-rank hypothesis. Consequently, that assessment does not establish the intended scope of all three questions. The original excerpt then runs into a separate Problem 12 concerning a Cartan--Hadamard question for median spaces, attributed there to Bowditch; that adjacent material is not part of the group-theoretic statement.
\subsection{Short English statement}
The supplied formulation asks whether every word-hyperbolic group acting properly and cocompactly on a finite-rank median space admits a geometric cubical action, is residually finite, and is linear. The original excerpt omits the hypotheses and linearity question, and the intended properness convention needs review. The supplied open status is a dated curation label, not an independently verified conclusion.
\subsection{Sources}
\begin{itemize}
\item Local UnsolvedMath registry, \texttt{data/unsolvedmath/problems.json}, permanent ID \texttt{30006623}, source code \texttt{catalogue entry}, title: \emph{Cubulability and Residual Finiteness of Hyperbolic Groups}. The supplied record belongs to the Oberwolfach Reports Open Problems collection and gives proposed year 2026; no proposer is identified.
\item Supplied report citation: \emph{Median Geometry and Applications} (2026), Oberwolfach Reports, pp.\ 487--540, DOI \texttt{10.4171/owr/2026/8}. Supplied source link: \url{https://doi.org/10.4171/owr/2026/8}.
\item Supplied literature source, labeled ``OWR report'': \url{https://ems.press/content/serial-article-files/53603}. The supplied assessment identifies Problem 11(2) as asking about residual finiteness, reports that the report presents it as open, and says that no resolution was found in that review.
\item Literature assessment date: \texttt{2026-08-22}. This is supplied dated curation, not independently verified here. The background explicitly calls it dated literature triage derived from the supplied OWR curation and requests independent verification of time-sensitive status and cited claims.
\item The registry labels the status \texttt{open} and the statement status \texttt{corrected\_verified}. Its verification note says that only formula delimiters, notation, grammar, and typography were repaired and that mathematical content and scope were preserved. These are supplied registry assertions; the discrepancies between the original excerpt and curated statement still require review.
\end{itemize}
\end{document}
